% ======================================================================
% Numerical experiments for the Affine-DDPNM paper
% Geometry-dependent modal requirements of finite-area interface
% representations in domain decomposition with pressure.
%
% Standalone compilation:
%   pdflatex numerical_experiments.tex
% (figures live in ./figures relative to this file)
%
% Section numbers 6-10 correspond to the intended position in the full
% paper (after the theory sections 1-5).  Only results computed with the
% benchmark pipelines archived under outputs/benchmark_w1n are reported;
% missing experiments are explicitly identified as pending.
% ======================================================================
\documentclass[11pt,a4paper]{article}

\usepackage[margin=2.4cm]{geometry}
\usepackage{amsmath,amssymb,amsthm}
\usepackage{bm}
\usepackage{booktabs}
\usepackage{graphicx}
\usepackage{xcolor}
\usepackage{hyperref}
\usepackage{enumitem}
\usepackage{array}

\newcommand{\R}{\mathbb{R}}
\newcommand{\dom}{\Omega}
\newcommand{\pore}[1]{\Omega_{#1}}
\newcommand{\iface}{\gamma}
\newcommand{\vel}{\bm{u}}
\newcommand{\pres}{p}
\newcommand{\normal}{\bm{n}}
\newcommand{\trac}{\bm{\lambda}}
\newcommand{\Schur}{\mathbf{S}}
\newcommand{\Shat}{\widehat{\Schur}}
\newcommand{\Shatuu}{\widehat{\Schur}{}^{uu}}
\newcommand{\Wspace}{\mathcal{W}}
\newcommand{\xiX}[1]{\bm{\xi}_{#1}}
\newcommand{\norm}[1]{\left\|#1\right\|}
\newcommand{\abs}[1]{\left|#1\right|}
\newcommand{\meas}[1]{\operatorname{meas}\left(#1\right)}
\newcommand{\cov}[1]{\operatorname{CV}\left(#1\right)}

\hypersetup{colorlinks=true,linkcolor=blue,urlcolor=blue,citecolor=blue}
\setlength{\emergencystretch}{2em}

\newtheorem{remark}{Remark}[section]

\begin{document}

\title{\textbf{Numerical Experiments}\\[4pt]
       \large Geometry-Dependent Modal Requirements of Finite-Area
       Interface Representations in Affine-DDPNM\\[4pt]
       \normalsize Standalone working version of \S6--\S10 of the full paper}
\author{}
\date{}
\maketitle

\setcounter{section}{5}

\noindent\textbf{Current numerical scope.}
Only completed benchmark results are assigned numerical values below.
The monolithic FEM reference and the theoretical broken-pressure C-DD
reference are kept distinct throughout.  Experiments that remain to be
run are identified explicitly in \S\ref{sec:limitations}; in particular,
no full-C-DD Schur-energy values or $\Wspace_{0v}$ accuracy values are
inferred from the available FEM-relative data.

% ======================================================================
\section{Numerical methodology}
\label{sec:methodology}
% ======================================================================

\subsection{Geometries and discretisations}
\label{sec:geometries}

Three three-dimensional test geometries are used.  All are unit-cube
packings of solid spheres with Stokes flow in the remaining fluid domain,
partitioned into $N_p$ pore subdomains with $m$ planar internal
interfaces.

\begin{enumerate}[label=(\roman*),nosep,leftmargin=*]
\item \textbf{Uniform-27}: a regular $3\times3\times3$ lattice of 27
  equal spheres (radius $0.105$) centred at
  $\{0.20,0.50,0.80\}^3$.  The partition is the regular $4\times4\times4$
  grid cut, giving $N_p = 64$ pores and $m = 144$ interfaces
  (mean $2m/N_p = 4.5$ ports per pore).
\item \textbf{Random-27}: a frozen random packing of 27 spheres
  (seed 20260804, 18 wall-clipped and 9 interior spheres); the partition
  is the Voronoi tessellation of the sphere centres, giving $N_p = 27$
  pores and $m = 114$ interfaces (mean 8.44 ports per pore).
\item \textbf{Heterogeneous-100}: a synthetic frozen random packing of 100 spheres
  (project-internal label ``Real-100'', seed 20260805) with porosity $0.856$;
  the partition is the regular
  $4\times4\times4$ grid cut, giving $N_p = 64$ pores and $m = 144$
  interfaces (mean 4.5 ports per pore).
\end{enumerate}

The solid-phase packings are shown schematically in
Fig.~\ref{fig:geometry} (the finite-element meshes themselves are not
rendered; the sphere positions and radii are those used to build the
meshes).  Mesh statistics are summarised in Table~\ref{tab:geometry},
including the objective variability descriptors suggested by the review
of the experimental design:
\begin{equation}
  \cov{A_\gamma} = \frac{\operatorname{std}(A_\gamma)}{\operatorname{mean}(A_\gamma)},
  \qquad
  \cov{V_i} = \frac{\operatorname{std}(V_i)}{\operatorname{mean}(V_i)},
  \qquad
  \cov{a_\gamma/b_\gamma} = \frac{\operatorname{std}(a_\gamma/b_\gamma)}
       {\operatorname{mean}(a_\gamma/b_\gamma)},
\end{equation}
where $A_\gamma$ is the area of interface $\gamma$, $V_i$ the volume of
pore $i$, and $a_\gamma/b_\gamma\ge 1$ the in-plane aspect ratio of the
interface patch, estimated as the square root of the ratio of the two
largest principal-component eigenvalues of the patch vertices.  These
are \emph{reported} as descriptive quantities; no claim is made that a
single scalar controls the error.

\begin{figure}[t]
  \centering
  \includegraphics[width=0.97\textwidth]{figures/fig1_geometry.png}
  \caption{Schematic of the three solid-phase packings used to build the
    test geometries (sphere positions and radii as used by the mesh
    generator; finite-element meshes are not shown): (a)~Uniform-27
    regular $3\times3\times3$ lattice; (b)~Random-27 frozen packing of
    27 spheres; (c)~Heterogeneous-100 frozen packing of 100 spheres.}
  \label{fig:geometry}
\end{figure}

\begin{table}[t]
  \centering
  \footnotesize
  \caption{Geometry and discretisation summary.  FEM DOF counts are for
    the monolithic Taylor--Hood pair $[P_2]^3$--$P_1$ (velocity DOFs =
    $3\times$ number of P2 nodes; pressure DOFs = number of P1 nodes);
    they match the mixed-unknown counts reported by the solver.}
  \label{tab:geometry}
  \resizebox{\textwidth}{!}{%
  \begin{tabular}{lccccccccl}
    \toprule
    Geometry & spheres & $N_p$ & $m$ & cells & vel. & pres. &
    mixed & $2m/N_p$ & remarks \\
    \midrule
    Uniform-27 & 27 & 64 & 144 & 12\,203 & 60\,936 & 3\,019 &
    63\,955 & 4.5 & regular lattice \\
    Random-27  & 27 & 27 & 114 & 15\,249 & 72\,471 & 3\,483 &
    75\,954 & 8.44 & Voronoi, irregular \\
    Heterogeneous-100   & 100 & 64 & 144 & 32\,611 & 159\,186 & 7\,655 &
    166\,841 & 4.5 & grid cut, porosity $0.856$ \\
    \bottomrule
  \end{tabular}}
  \medskip
  \begin{tabular}{lcccc}
    \toprule
    Geometry & $\cov{V_i}$ & $\cov{A_\gamma}$ & $\cov{a_\gamma/b_\gamma}$ &
    mean $a_\gamma/b_\gamma$ \\
    \midrule
    Uniform-27 & 0.321 & 0.203 & 0.158 & 1.39 \\
    Random-27  & 0.432 & 0.783 & 1.479 & 2.48 \\
    Heterogeneous-100   & 0.057 & 0.116 & 0.078 & 1.11 \\
    \bottomrule
  \end{tabular}
\end{table}

The three configurations comprise a deliberately regular lattice
(Uniform-27), a deliberately irregular Voronoi-partitioned case
(Random-27), and a larger synthetic heterogeneous packing
(Heterogeneous-100).  The partition-level descriptors in
Table~\ref{tab:geometry} should not be interpreted as a scalar measure of
``geometric complexity'': Heterogeneous-100 uses a regular grid partition
and therefore has small $\cov{V_i}$, $\cov{A_\gamma}$ and
$\cov{a_\gamma/b_\gamma}$ despite its heterogeneous solid arrangement.

\subsection{Reference systems and error metrics}
\label{sec:references}

Two reference systems must be kept distinct (see also the theory
document):

\begin{enumerate}[label=(\roman*),nosep,leftmargin=*]
\item \textbf{Physical/numerical reference}: the monolithic Taylor--Hood
  finite-element solution
  $\vel_h^{\mathrm{FEM}}, \pres_h^{\mathrm{FEM}}$ on the \emph{same}
  mesh, with continuous pressure across pores.  All reported accuracy
  metrics are relative to this reference.  The theoretical section does
  not establish equivalence between this system and the broken-pressure
  C-DD reference; the error relative to the FEM solution is therefore
  reported as an \emph{observed} accuracy measure, and no rate statement
  is made about it from the theory.
\item \textbf{Theoretical Galerkin reference}: the broken-pressure C-DD
  trace solution $\xiX{\mathrm{CDD}}$ of the theory sections, and the
  Schur-energy seminorm $\norm{\cdot}_{\Shatuu}$ it induces.  The
  theory-level quantities ($\norm{\xiX{\mathrm{CDD}}-\xiX{X}}_{\Shatuu}$,
  the Pythagorean decompositions) were \emph{not} computed in the runs
  archived so far (the full C-DD solve was not executed); they are
  discussed in \S\ref{sec:theoryverify} as planned verification, and no numerical values
  are reported for them.
\end{enumerate}

For each reduced space $X \in \{\Wspace_{0n}, \Wspace_{0v},
\Wspace_{1n}, \Wspace_{1v}\}$ we report, uniformly:
\begin{align}
  e_u^{L^2} &= \frac{\norm{\vel_X - \vel_{\mathrm{FEM}}}_{L^2(\dom)}}
                   {\norm{\vel_{\mathrm{FEM}}}_{L^2(\dom)}}, \label{eq:elu}\\
  e_u^{H^1_{\mathrm{br}}} &= \frac{
      \left(\sum_i \norm{\nabla(\vel_X-\vel_{\mathrm{FEM}})}_{L^2(\pore{i})}^2\right)^{1/2}}
      {\left(\sum_i \norm{\nabla \vel_{\mathrm{FEM}}}_{L^2(\pore{i})}^2\right)^{1/2}},
      \label{eq:ehu}\\
  e_p^{L^2} &= \frac{\norm{\pres_X - \pres_{\mathrm{FEM}}}_{L^2(\dom)}}
                   {\norm{\pres_{\mathrm{FEM}}}_{L^2(\dom)}}, \label{eq:elp}\\
  e_Q &= \frac{\abs{Q_X - Q_{\mathrm{FEM}}}}{\abs{Q_{\mathrm{FEM}}}},
       \qquad Q = \text{outlet volume flux}. \label{eq:eq}
\end{align}
All error norms are evaluated cell-wise on the common mesh (broken
domain), with degree-6 quadrature.  The pressure gauge is fixed by
Dirichlet pressure conditions at inlet ($p=1$) and outlet ($p=0$) in
both the FEM and the reduced systems; the raw relative $L^2$ pressure
error \eqref{eq:elp} is reported without any post-hoc mean alignment.
Because the pressure level is fixed physically by the inlet/outlet
conditions, this is the pressure metric used consistently below.  The
theory gives no monotonicity of \eqref{eq:elp} along the nested chains,
so occasional non-monotone pressure errors are not a discrepancy.

\subsection{Four modal interface spaces}
\label{sec:spaces}

The numerical programme is organised around the four trace subspaces
of the theory section, whose per-interface contents are
\begin{align}
  \Wspace_{0n} &: \{1\cdot \normal\}, \qquad
  \Wspace_{0v} : \{1\cdot \normal,\; 1\cdot \bm{t}_1,\; 1\cdot \bm{t}_2\},\\
  \Wspace_{1n} &: \{1\cdot \normal,\; s\cdot \normal,\; t\cdot \normal\},\\
  \Wspace_{1v} &: \{1,s,t\}\otimes\{\normal,\bm{t}_1,\bm{t}_2\},
\end{align}
where $(s,t)$ are local planar coordinates on the (representative)
interface and $\{\normal,\bm{t}_1,\bm{t}_2\}$ a fixed orthonormal frame.
Table~\ref{tab:spaces} lists the unknown budgets.  The two three-mode
spaces satisfy $\dim\Wspace_{0v}=\dim\Wspace_{1n}=3m$, which makes them
an \emph{equal-budget} pair: the same number of interface unknowns spent
on directional content ($\Wspace_{0v}$) versus on first-order spatial
content ($\Wspace_{1n}$).  Of the two paths
$\Wspace_{0n}\subseteq\Wspace_{0v}\subseteq\Wspace_{1v}$ and
$\Wspace_{0n}\subseteq\Wspace_{1n}\subseteq\Wspace_{1v}$, only the
second was executable with the implemented bases at the time of the
archived runs; $\Wspace_{0v}$ (constant tangential modes) is listed as
\emph{not computed} in this study.

\begin{table}[t]
  \centering
  \small
  \caption{The four modal interface spaces and their global unknown
    budgets on the three geometries ($m$ = number of interfaces:
    144/114/144).  The row marked $\dagger$ was \emph{not computed} in
    the runs archived so far.}
  \label{tab:spaces}
  \begin{tabular}{lllccc}
    \toprule
    Method & space/interface & modes/iface & \multicolumn{3}{c}{global unknowns}\\
    \cmidrule(lr){4-6}
      & & & Uniform-27 & Random-27 & Heterogeneous-100 \\
    \midrule
    Classic       & $\Wspace_{0n}$ & 1 & 144 & 114 & 144 \\
    P0-vector\dag & $\Wspace_{0v}$ & 3 & 432 & 342 & 432 \\
    NormalLinear  & $\Wspace_{1n}$ & 3 & 432 & 342 & 432 \\
    Affine        & $\Wspace_{1v}$ & 9 & 1296 & 1026 & 1296 \\
    \bottomrule
  \end{tabular}
  \medskip
  \raggedright\footnotesize
  Per-interface contents: $\Wspace_{0n}=\{1\normal\}$;
  $\Wspace_{0v}=\{1\normal,1\bm{t}_1,1\bm{t}_2\}$;
  $\Wspace_{1n}=\{1\normal,s\normal,t\normal\}$;
  $\Wspace_{1v}=\{1,s,t\}\otimes\{\normal,\bm{t}_1,\bm{t}_2\}$.
  $\dag$ Constant-vector (P0-vector) space: basis implemented as part
  of the four-space framework but not run in the archived benchmarks;
  included for the equal-budget comparison in \S\ref{sec:equalbudget}.
\end{table}

\subsection{Timing and solver protocol}
\label{sec:timing}

All archived runs used the same serial pipeline on one Windows 11
x86-64 workstation, with viscosity $1$ and inlet/outlet pressure $1/0$.
The exact CPU/RAM and software-version metadata are not recorded in this
working draft and must be added for the final performance benchmark.  Timings distinguish
\begin{align}
  T_{\mathrm{offline}} &= \text{local factorisations and primitive
    response construction},\\
  T_{\mathrm{online}} &= \text{global reduced assembly + Schur solve +
    local reconstruction}.
\end{align}
$T_{\mathrm{first}} = T_{\mathrm{offline}} + T_{\mathrm{online}}$ is the
wall time of a first reduced solve.  Common mesh generation and
post-processing validation are excluded.  The archived timings are not
yet a fully uniform performance benchmark: for Uniform-27 and Random-27
the FEM entry is assembly+solve, whereas the Heterogeneous-100 archive
contains FEM solve time only.  Consequently Table~\ref{tab:cost} reports
absolute timing scopes explicitly.  First-solve speedups are quoted only
where the FEM and reduced timings are directly comparable; online-only
ratios for Heterogeneous-100 are labelled as such and are not interpreted
as single-solve wall-clock speedups.

% ======================================================================
\section{Verification of the reference implementation}
\label{sec:verification}
% ======================================================================

\subsection{Exact-FE-Schur: Schur elimination code validation}
\label{sec:exactschur}

On Random-27 the benchmark additionally solves the exact Schur
complement of the monolithic system on the full FE interface trace
(21\,970 trace DOFs, 465.9\,s).  This serves as a code-level validation
of the Schur-elimination machinery, \emph{not} as the theoretical C-DD
reference: the two systems share the assembled monolithic operator and
must agree to round-off.  The measured relative difference between the
exact-FE-Schur solution and the monolithic solve is
\begin{equation}
  \frac{\norm{\vel_{\mathrm{Schur}}-\vel_{\mathrm{direct}}}}
       {\norm{\vel_{\mathrm{direct}}}} = 1.05\times10^{-12},
\end{equation}
with Schur symmetry defect $3.2\times10^{-15}$, Schur relative residual
$5.75\times10^{-14}$ and global relative residual $1.55\times10^{-15}$.
This validates the full-trace Schur elimination and the
reference/error-evaluation pipeline.  The reduced systems are checked
separately through their Schur residuals, symmetry and conservation
diagnostics below; the Exact-FE-Schur comparison is not, by itself, a
validation of the affine modal-basis assembly.

\subsection{Conservation and algebraic diagnostics}
\label{sec:algebraic}

Table~\ref{tab:diagnostics} collects the algebraic diagnostics recorded
by the runs.  Local discrete incompressibility gives each pore balance,
while the retained constant-normal Schur equation enforces cancellation
of the internal interface fluxes; together these yield the global mass
balance under the assumptions of the conservation theorem.  The
residuals shown are the maximum absolute modal-moment residual entries of
the assembled Schur systems, at round-off level in all recorded cases.  The Schur matrices are
exactly symmetric by construction (symmetry defect zero for the reduced
systems).  On Random-27 the recorded minimum Schur eigenvalue drops from
$2.1\times10^{-8}$ (Classic) to $1.26\times10^{-12}$ for both
NormalLinear and Affine.  This indicates a much smaller minimum spectral
scale in the chosen coefficient normalization, but it is not by itself a
condition-number estimate because $\lambda_{\max}(\Schur)$ was not
archived.  The matrices were successfully factorised as SPD in these
runs.

\begin{table}[t]
  \centering
  \small
  \caption{Numerical verification diagnostics as recorded by the
    archived runs (``--'' = not recorded by that benchmark).
    ``$\max r_{\mathrm{mom}}$'' is the maximum absolute modal-moment
    residual of the assembled reduced system; the mass-balance row is
    the inlet/outlet flux sum of the random benchmark.}
  \label{tab:diagnostics}
  \begin{tabular}{lccccc}
    \toprule
    Geometry & Method & $\max\abs{r_{\mathrm{mom}}}$ & Schur sym.\ defect &
    $\lambda_{\min}(\Schur)$ & in/out flux sum \\
    \midrule
    Uniform-27 & all & -- & -- & $1.67\times10^{-4}$\dag & -- \\
    \addlinespace[2pt]
    Random-27 & $W_{0n}$ & $2.17\times10^{-16}$ & 0 & $2.11\times10^{-8}$ &
      $-5\times10^{-16}$ \\
    Random-27 & $W_{1n}$ & $3.55\times10^{-16}$ & 0 & $1.26\times10^{-12}$ &
      $-2\times10^{-16}$ \\
    Random-27 & $W_{1v}$ & $4.92\times10^{-16}$ & 0 & $1.26\times10^{-12}$ &
      $-4\times10^{-16}$ \\
    \addlinespace[2pt]
    Heterogeneous-100 & $W_{0n}$ & $6.36\times10^{-17}$ & 0 & -- & -- \\
    Heterogeneous-100 & $W_{1n}$ & $1.05\times10^{-16}$ & 0 & -- & -- \\
    Heterogeneous-100 & $W_{1v}$ & $1.45\times10^{-16}$ & 0 & -- & -- \\
    \bottomrule
  \end{tabular}
  \medskip
  \raggedright\footnotesize
  $\dag$ Recorded by the base-project verification run
  (\texttt{ddpnm\_3d\_uniform\_spheres}, same geometry, Classic system);
  the affine benchmark did not record the eigenvalue.  The monolithic
  FEM reference solves reached relative linear residuals
  $1.73\times10^{-12}$ (Uniform-27) and $2.92\times10^{-12}$ (Random-27),
  with relative mass imbalance $2.54\times10^{-14}$ (Random-27).
\end{table}

\subsection{Theory-level verification: not yet computed}
\label{sec:theoryverify}

The theory sections derive the best-approximation property in the
$\Shatuu$-seminorm and the exact Pythagorean decompositions along the
two nested paths, with the reduction errors
$\Delta_T=\norm{\xiX{0v}-\xiX{0n}}_{\Shatuu}^2$,
$\Delta_N=\norm{\xiX{1n}-\xiX{0n}}_{\Shatuu}^2$, and the identity
\begin{equation}
  \norm{\xiX{\mathrm{CDD}}-\xiX{0n}}_{\Shatuu}^2 =
  \norm{\xiX{\mathrm{CDD}}-\xiX{1v}}_{\Shatuu}^2
  + \norm{\xiX{1v}-\xiX{1n}}_{\Shatuu}^2
  + \norm{\xiX{1n}-\xiX{0n}}_{\Shatuu}^2,
\end{equation}
and its analogue along $\Wspace_{0n}\to\Wspace_{0v}\to\Wspace_{1v}$.
These identities are exact for nested spaces in exact arithmetic and
their numerical verification requires solving the \emph{full} C-DD
system (the broken-pressure Galerkin solution $\xiX{\mathrm{CDD}}$),
which was not executed in the archived runs.  This verification is a
planned item (see \S\ref{sec:limitations}).  The nestedness of the
implemented bases is exact by construction, while
\S\ref{sec:exactschur}--\S\ref{sec:algebraic} verify the full-trace
Schur elimination and the algebraic residual/conservation properties of
the reduced systems.  They do not establish equivalence between a
reduced operator and the monolithic FEM operator.

% ======================================================================
\section{Controlled modal experiments}
\label{sec:modal}
% ======================================================================

\subsection{Absolute accuracy}
\label{sec:accuracy}

Table~\ref{tab:master} collects the complete accuracy results of all
archived runs: 3 geometries $\times$ 3 computed spaces, four metrics
each.  It is an absolute-results table (not improvement ratios).

\begin{table}[t]
  \centering
  \small
  \caption{Complete accuracy results relative to the monolithic
    Taylor--Hood FEM on the same mesh.  $e_u^{L^2}$,
    $e_u^{H^1_{\mathrm{br}}}$, $e_p^{L^2}$ and $e_Q$ are defined in
    \eqref{eq:elu}--\eqref{eq:eq}.  The $\Wspace_{0v}$ row was not
    computed in the archived runs.}
  \label{tab:master}
  \begin{tabular}{llcccc}
    \toprule
    Geometry & Method & $e_u^{L^2}$ & $e_u^{H^1_{\mathrm{br}}}$ &
    $e_p^{L^2}$ & $e_Q$ \\
    \midrule
    Uniform-27 & $W_{0n}$ & 0.2108 & 0.4483 & 0.0321 & 0.1646 \\
    Uniform-27 & $W_{1n}$ & 0.0816 & 0.2385 & 0.0185 & 0.0243 \\
    Uniform-27 & $W_{1v}$ & 0.0563 & 0.2161 & 0.0217 & 0.0032 \\
    \addlinespace[2pt]
    Random-27 & $W_{0n}$ & 0.6532 & 0.8585 & 0.1030 & 0.7272 \\
    Random-27 & $W_{1n}$ & 0.3075 & 0.3897 & 0.0322 & 0.1465 \\
    Random-27 & $W_{1v}$ & 0.0651 & 0.1756 & 0.0203 & 0.0271 \\
    \addlinespace[2pt]
    Heterogeneous-100 & $W_{0n}$ & 0.9254 & 1.0627 & 0.0776 & 1.1199 \\
    Heterogeneous-100 & $W_{1n}$ & 0.4738 & 0.5129 & 0.0299 & 0.2595 \\
    Heterogeneous-100 & $W_{1v}$ & 0.1998 & 0.3594 & 0.0236 & 0.1244 \\
    \bottomrule
  \end{tabular}
\end{table}

Three observations follow directly from the absolute numbers, all
qualitative and confined to the tested configurations:

\begin{enumerate}[label=\textbf{(C\arabic*)}.,leftmargin=*,nosep]
\item The Classic constant-normal model $\Wspace_{0n}$ exhibits large
  FEM-relative velocity errors on all three tested configurations
  ($21.1\%$, $65.3\%$ and $92.5\%$, respectively).  These are observed
  model-to-FEM discrepancies; without the full C-DD reference they are
  not identified purely as modal-reduction errors.
\item Along the computed path $\Wspace_{0n}\to\Wspace_{1n}\to
  \Wspace_{1v}$, the first-order normal step accounts for the majority
  of the observed Classic-to-Affine $L^2$ velocity-error reduction on
  all three configurations (quantified in \S\ref{sec:shares}).
\item The additional full-vector affine enrichment remains substantial
  for Random-27 ($e_u^{L^2}$: $0.308\to0.065$) and Heterogeneous-100
  ($0.474\to0.200$), while it is a smaller field-error correction on
  Uniform-27 ($0.082\to0.056$).  On Uniform-27 the outlet-flux error
  drops to $0.32\%$ only with $\Wspace_{1v}$.
\end{enumerate}

The non-monotonicity of the pressure error on Uniform-27
($e_p^{L^2}(\Wspace_{1n}) = 0.0185 < e_p^{L^2}(\Wspace_{1v}) = 0.0217$)
is not a discrepancy: the theory does not guarantee monotonicity of the
pressure $L^2$ error along nested chains (the reduction error is
monotone in the Schur-energy seminorm, not in $L^2(p)$).

\subsection{Accuracy versus modal budget}
\label{sec:budget}

Fig.~\ref{fig:budget} is the central visual of the section: for each
geometry it plots the relative $L^2$ velocity error against the number
of modes per interface, $r \in \{1,3,9\}$, joining the points along the
computed nested path $\Wspace_{0n}\to\Wspace_{1n}\to\Wspace_{1v}$.  The
$r=3$ location carries a single computed space ($\Wspace_{1n}$); the
equal-budget companion $\Wspace_{0v}$ is not yet computed, so the
question ``same DOF budget spent on direction versus on spatial
variation'' can only be posed here, not answered (see \S\ref{sec:equalbudget}).

\begin{figure}[t]
  \centering
  \includegraphics[width=0.97\textwidth]{figures/fig2_modal_budget.png}
  \caption{Relative $L^2$ velocity error versus modes per interface
    ($r=1,3,9$) along the nested path
    $\Wspace_{0n}\to\Wspace_{1n}\to\Wspace_{1v}$, for the three
    geometries (log scale).  Only the path through $\Wspace_{1n}$ was
    computed; the equal-budget space $\Wspace_{0v}$ at $r=3$ is
    planned work.}
  \label{fig:budget}
\end{figure}

In absolute FEM-relative $L^2$ error, the reduction from $r=1$ to
$r=3$ is larger than the reduction from $r=3$ to $r=9$ in each of the
three archived cases.  This statement is intentionally additive rather
than multiplicative: the log-scale ratios need not exhibit the same
ordering.  The remaining error at $r=9$ also differs substantially
($5.6\%$, $6.5\%$, $20.0\%$), indicating configuration-dependent
accuracy at a fixed modal budget; no asymptotic statement in $r$ is
claimed.

\subsection{Metric comparison across geometries}
\label{sec:metrics}

Fig.~\ref{fig:metrics} shows all four metrics as dot plots
(velocity $L^2$, broken $H^1$, pressure $L^2$, outlet-flux error) with
logarithmic ordinate; the wide dynamic range ($0.3\%$--$112\%$) is the
reason for the log scale.  The figure reproduces the Table~\ref{tab:master} numbers in
a form where the cross-geometry pattern is visible at a glance: the
Classic-to-Affine spread widens from the regular to the irregular and
heterogeneous cases in every metric except the pressure error, which stays
within one decade across all geometries and methods.

\begin{figure}[t]
  \centering
  \includegraphics[width=0.97\textwidth]{figures/fig3_metrics.png}
  \caption{Four error metrics (relative, log scale) across the three
    geometries, three computed spaces: (a)~velocity $L^2$,
    (b)~velocity broken $H^1$, (c)~pressure $L^2$, (d)~outlet-flux
    error.}
  \label{fig:metrics}
\end{figure}

\subsection{Equal-budget comparison \texorpdfstring{$\Wspace_{0v}$ vs.\ $\Wspace_{1n}$}{W0v vs. W1n}: planned}
\label{sec:equalbudget}

The central controlled contrast of the programme is the equal-budget
comparison
\begin{equation}
  G_{0v} = \frac{e_{0n}-e_{0v}}{e_{0n}}, \qquad
  G_{1n} = \frac{e_{0n}-e_{1n}}{e_{0n}},
\end{equation}
i.e.\ the FEM-relative $L^2$ reduction obtained by spending three DOFs
per interface on constant directional content ($\Wspace_{0v}$) versus
on first-order spatial content ($\Wspace_{1n}$).  This experiment was
\emph{not} run in the archived benchmarks (the constant-vector basis is
implemented in the four-space framework but no runs were executed).
It is the first item of \S\ref{sec:limitations}.  Until it is computed, the statement
that the normal-linear modes are the dominant mechanism rests on the
comparison $\Wspace_{0n}\to\Wspace_{1n}\to\Wspace_{1v}$ alone
(Fig.~\ref{fig:budget}), which cannot separate directional from spatial
information.

\subsection{Field visualization (Random-27)}
\label{sec:fields}

Fig.~\ref{fig:fields} shows, on the Random-27 geometry, a slice plane
($z=0.5$) of the FEM velocity magnitude together with the pointwise
velocity error magnitude of the three reduced spaces, using the same
error-map color normalization.  Qualitatively, the error becomes both
weaker and more spatially localized along the sequence
$\Wspace_{0n}\to\Wspace_{1n}\to\Wspace_{1v}$, consistent with the
global metrics.  Because the archived figure does not retain numerical
colorbar limits, it is used only as a qualitative field visualization;
all quantitative comparisons use Tables~\ref{tab:master} and
\ref{tab:cross}.

\begin{figure}[t]
  \centering
  \includegraphics[width=0.97\textwidth]{figures/fig6_fields.png}
  \caption{Random-27, slice $z=0.5$: FEM velocity magnitude (left) and
    pointwise velocity-error magnitude of $\Wspace_{0n}$,
    $\Wspace_{1n}$ and $\Wspace_{1v}$ with a common archived color
    normalization.  The panel is qualitative because numerical colorbar
    limits were not retained in the archived image.}
  \label{fig:fields}
\end{figure}

The archived slice output is available only for Random-27.  Additional
field visualizations for the other configurations remain to be generated
from the corresponding solution fields.

% ======================================================================
\section{Larger heterogeneous configuration (Heterogeneous-100)}
\label{sec:realistic}
% ======================================================================

\subsection{Accuracy}
\label{sec:realacc}

Heterogeneous-100 is the largest of the three synthetic configurations
(32\,611 cells, 166\,841 mixed DOFs, porosity $0.856$).  Its accuracy
rows in Table~\ref{tab:master} repeat the qualitative pattern of the
controlled geometries, at a higher level: Classic $\Wspace_{0n}$ has a $92.5\%$ relative velocity
$L^2$ error and a broken-$H^1$ relative error above one ($e_u^{H^1_{\mathrm{br}}} = 1.063$); NormalLinear
reduces the $L^2$ error to $47.4\%$ and the flux error to $25.9\%$;
Affine reaches $e_u^{L^2}=20.0\%$, $e_Q=12.4\%$.  The pressure error is
the least sensitive metric here as well ($7.8\%\to3.0\%\to2.4\%$).
These are the largest errors among the three tested configurations for
each of the three computed spaces.  The present data do not isolate a
single geometric or topological quantity responsible for this ordering.

\subsection{Cost and memory}
\label{sec:realcost}

Table~\ref{tab:cost} reports the cost structure of all runs.  For
Uniform-27 and Random-27 the reduced methods achieve a genuine
first-solve speedup over the monolithic FEM ($10.5\times$ and $13.3\times$
for Classic; $8.2\times$ and $9.1\times$ for Affine), because the
offline library is cheap relative to the FEM solve on those meshes.
For Heterogeneous-100 the archive contains a $14.5\,$s FEM solve-only
time, while the reduced offline phase takes $28$--$33\,$s.  Thus the
archived data do not support a like-for-like first-solve speedup claim
for this case.  The reduced online phase is short ($0.06$--$1.4\,$s),
and the quoted $255\times$, $72\times$ and $10\times$ ratios compare
that online phase against the archived FEM solve time only.  A formal
performance table should be rerun with one timing scope and repeated
measurements on all configurations.

Memory: peak footprints on Heterogeneous-100 are $45$\,MiB (Classic),
$79$\,MiB (NormalLinear), $246$\,MiB (Affine) versus $300$\,MiB for the
monolithic solve.  The observed increase with mode count is consistent
with the larger reduced coupling blocks, but no memory-complexity law is
inferred from these three measurements.

\begin{table}[t]
  \centering
  \small
  \caption{Computational cost (seconds).  $T_{\mathrm{offline}}$: local
    factorisations + primitive responses; $T_{\mathrm{online}}$:
    reduced assembly + Schur solve + reconstruction;
    $T_{\mathrm{first}} = T_{\mathrm{offline}}+T_{\mathrm{online}}$.
    $N_{\mathrm{prim}}$ = primitive RHS columns.  For Uniform-27 and
    Random-27 the FEM time is assembly+solve; for Heterogeneous-100 the
    archived FEM entry is solve-only, so timing ratios for that block are
    labelled ``online only'' and are not directly comparable to the
    first-solve ratios above.}
  \label{tab:cost}
  \resizebox{\textwidth}{!}{%
  \begin{tabular}{llcccccc}
    \toprule
    Geometry & Method & unknowns & $N_{\mathrm{prim}}$ &
    $T_{\mathrm{offline}}$ & $T_{\mathrm{online}}$ &
    $T_{\mathrm{first}}$ & speedup vs FEM \\
    \midrule
    Uniform-27 & $W_{0n}$ & 144 & 320 & 7.35 & 0.0071 & 7.36 & $10.50\times$ \\
    Uniform-27 & $W_{1n}$ & 432 & 896 & 7.38 & 0.0222 & 7.40 & $10.44\times$ \\
    Uniform-27 & $W_{1v}$ & 1296 & 2624 & 9.18 & 0.208 & 9.39 & $8.23\times$ \\
    Uniform-27 & FEM & 63\,955 & -- & -- & 77.25 & 77.25 & $1\times$ \\
    \addlinespace[2pt]
    Random-27 & $W_{0n}$ & 114 & 246 & 8.90 & 0.0141 & 8.91 & $13.29\times$ \\
    Random-27 & $W_{1n}$ & 342 & 702 & 9.20 & 0.0287 & 9.22 & $12.84\times$ \\
    Random-27 & $W_{1v}$ & 1026 & 2070 & 12.82 & 0.191 & 13.01 & $9.10\times$ \\
    Random-27 & FEM & 75\,954 & -- & -- & 118.42 & 118.42 & $1\times$ \\
    Random-27 & Exact-FE-Schur & 21\,970 & -- & -- & 465.90 & 465.90 & $0.25\times$ \\
    \addlinespace[2pt]
    Heterogeneous-100 & $W_{0n}$ & 144 & -- & 27.9 & 0.057 & 27.98 & $255\times$ (online only) \\
    Heterogeneous-100 & $W_{1n}$ & 432 & -- & 28.0 & 0.203 & 28.17 & $72\times$ (online only) \\
    Heterogeneous-100 & $W_{1v}$ & 1296 & -- & 32.7 & 1.400 & 34.13 & $10\times$ (online only) \\
    Heterogeneous-100 & FEM & 166\,841 & -- & -- & 14.53 & 14.53 & $1\times$ \\
    \bottomrule
  \end{tabular}}
  \medskip
  \raggedright\footnotesize
  Peak memory (Heterogeneous-100 only): $W_{0n}$ 45.0, $W_{1n}$ 78.8,
  $W_{1v}$ 246.4, FEM 299.8\,MiB.  ``Online only'' speedups are
  $T_{\mathrm{online}}(\mathrm{FEM})/T_{\mathrm{online}}(X)$; these
  ratios exclude the reduced offline library and therefore do
  not constitute single-solve speedups.
\end{table}

\subsection{Accuracy--cost tradeoff}
\label{sec:pareto}

Fig.~\ref{fig:pareto} plots the relative $L^2$ velocity error of the
three reduced spaces against their online time.  The monolithic FEM
solution has zero error relative to itself and therefore cannot be
plotted on a logarithmic error axis; its archived time is shown only as
a vertical reference line in each panel.  The figure visualizes the
reduced-method accuracy--online-cost tradeoff and is relevant primarily
to repeated-query use, because it excludes the offline library.  The
Heterogeneous-100 FEM timing scope also differs from the other two cases,
as stated in Table~\ref{tab:cost}.

\begin{figure}[t]
  \centering
  \includegraphics[width=0.97\textwidth]{figures/fig5_pareto.png}
  \caption{Relative $L^2$ velocity error versus online solve time
    (log--log), shown separately for the three configurations.  The
    vertical dotted line is the archived monolithic FEM time; the FEM
    self-error is exactly zero and is therefore off the logarithmic
    ordinate.  Online time excludes the reduced offline library (see
    Table~\ref{tab:cost}).}
  \label{fig:pareto}
\end{figure}

% ======================================================================
\section{Cross-geometry modal analysis}
\label{sec:cross}
% ======================================================================

\subsection{Observed FEM-relative error-reduction shares}
\label{sec:shares}

Along the computed path $\Wspace_{0n}\to\Wspace_{1n}\to\Wspace_{1v}$ the
observed FEM-relative $L^2$ velocity-error reduction decomposes into the
normal-spatial step ($\Wspace_{0n}\to\Wspace_{1n}$) and the additional
vectorial step ($\Wspace_{1n}\to\Wspace_{1v}$):
\begin{equation}
  R_N = \frac{e_{0n}-e_{1n}}{e_{0n}-e_{1v}}, \qquad
  R_{V|N} = \frac{e_{1n}-e_{1v}}{e_{0n}-e_{1v}}, \qquad
  R_N + R_{V|N} = 1.
  \label{eq:shares}
\end{equation}
Fig.~\ref{fig:shares} and Table~\ref{tab:cross} report these shares.

\begin{figure}[t]
  \centering
  \includegraphics[width=0.72\textwidth]{figures/fig4_shares.png}
  \caption{Fractions of the \emph{observed} FEM-relative $L^2$ velocity
    error reduction along the prescribed path
    $\Wspace_{0n}\to\Wspace_{1n}\to\Wspace_{1v}$ (100\% stacked).
    These fractions are descriptive; they are not unique causal modal
    contributions, because the two middle spaces $\Wspace_{0v}$ and
    $\Wspace_{1n}$ are generally non-nested and the path choice fixes the attribution.}
  \label{fig:shares}
\end{figure}

\begin{table}[t]
  \centering
  \small
  \caption{Cross-geometry modal summary: absolute errors and the
    observed FEM-relative reduction shares of \eqref{eq:shares}.}
  \label{tab:cross}
  \begin{tabular}{lcccccc}
    \toprule
    Geometry & $e_{0n}^{L^2}$ & $e_{1n}^{L^2}$ & $e_{1v}^{L^2}$ &
    $R_N$ & $R_{V|N}$ & reduction $\frac{e_{0n}-e_{1v}}{e_{0n}}$ \\
    \midrule
    Uniform-27 & 0.2108 & 0.0816 & 0.0563 & 83.6\% & 16.4\% & 73.3\% \\
    Random-27  & 0.6532 & 0.3075 & 0.0651 & 58.8\% & 41.2\% & 90.0\% \\
    Heterogeneous-100   & 0.9254 & 0.4738 & 0.1998 & 62.2\% & 37.8\% & 78.4\% \\
    \bottomrule
  \end{tabular}
\end{table}

Two statements can be made from these numbers, both confined to the
tested configurations and to the chosen path:

\begin{enumerate}[label=\textbf{(D\arabic*)}.,leftmargin=*,nosep]
\item On all three geometries the normal-spatial step
  $\Wspace_{0n}\to\Wspace_{1n}$ accounts for the majority of the
  observed reduction ($59$--$84\%$).
\item The share of the additional full-vector affine step
  $\Wspace_{1n}\to\Wspace_{1v}$ is larger for Random-27 and
  Heterogeneous-100 ($41\%$ and $38\%$) than for Uniform-27 ($16\%$).
  The simple partition-level variability descriptors in
  Table~\ref{tab:geometry} do \emph{not} order these three cases
  monotonically---Heterogeneous-100 has smaller listed CVs than
  Uniform-27---so the present data support configuration-dependent modal
  behavior but not a scalar ``irregularity $\mapsto$ modal share'' law.
\end{enumerate}

\subsection{Schur-energy pathwise decomposition: planned}
\label{sec:schurdecomp}

The theory sections give an \emph{exact} decomposition of the squared
Schur-energy error into pathwise marginal gains (see \S\ref{sec:theoryverify}).  Unlike
the FEM-relative shares of Fig.~\ref{fig:shares}, those quantities
($\Delta_N$, $\Delta_{T|N}$ etc.\ in the $\Shatuu$-seminorm) are
theorem-level and require the full C-DD trace solution
$\xiX{\mathrm{CDD}}$.  They were not computed in the archived runs and
are therefore not reported numerically; Fig.~\ref{fig:shares} and the
planned Schur-energy figure must not be conflated.  The former is an
observed descriptive attribution, the latter would be a strict theorem
quantity; only the former is available at this stage.

\subsection{Geometry dependence and limitations}
\label{sec:limitations}

\textbf{What the current data supports.}  (i)~the Classic
constant-normal model $\Wspace_{0n}$ exhibits large FEM-relative
velocity errors in all three tested configurations; (ii)~adding
first-order normal spatial content $\Wspace_{1n}$ produces a major
observed correction in each case; (iii)~the residual gap between
$\Wspace_{1n}$ and $\Wspace_{1v}$ is small for Uniform-27 and
substantial for Random-27 and Heterogeneous-100, indicating
configuration-dependent modal requirements across the tested cases;
(iv)~neither a monotonic dependence on a single scalar geometry
descriptor nor a causal geometric law is claimed.

\textbf{What is planned and not yet computed.}  (a)~the equal-budget
space $\Wspace_{0v}$ (constant tangential modes) on all three
geometries, enabling the controlled contrast of \S\ref{sec:equalbudget};
(b)~the full C-DD solve $\xiX{\mathrm{CDD}}$ on Uniform-27 and
Random-27 for the exact Schur-energy verification (\S\ref{sec:theoryverify}) and the
pathwise decomposition (\S\ref{sec:schurdecomp}); (c)~a mesh-sensitivity study
(Uniform-27 at three mesh levels) to confirm that the separation
$\Wspace_{0n}\gg\Wspace_{1n}>\Wspace_{1v}$ is not an artifact of one
mesh; the theory gives fixed-mesh separation only, no
mesh-independent asymptotic statement; (d)~an ensemble of random
realisations of the Random-27 geometry to test the stability of the
$59/41$ split across seeds; (e)~additional field outputs for
Heterogeneous-100; (f)~an independent second porous configuration at a
different pore/sphere count; and (g)~a uniform, repeated timing benchmark
with the same FEM/reduced timing scope on every configuration.

\textbf{Scope notes.}  All timings are single-machine, single-threaded
measurements from the archived runs; they are indicative, not
benchmark-suite grade.  The reported errors are with respect to the
monolithic Taylor--Hood FEM on the same mesh; the relation between the
broken-pressure C-DD discretisation and the monolithic system is not
established in the theory sections, and therefore no claim about the
accuracy of the reduced methods \emph{relative to the exact Stokes
solution} is made.

% ======================================================================
\section*{Appendix A: archived output provenance}
\label{app:provenance}
% ======================================================================

All numerical values in this working section were read from the archived
benchmark outputs listed below:
\begin{itemize}[nosep,leftmargin=*]
\item \texttt{affine\_ddpnm\_3d/outputs/benchmark\_w1n/}:\\
  \texttt{affine\_ddpnm\_metrics.csv}, \texttt{affine\_ddpnm\_report.json};
\item \texttt{affine\_ddpnm\_3d\_random\_porous/outputs/benchmark\_w1n/}:\\
  \texttt{random\_affine\_metrics.csv}, \texttt{random\_affine\_report.json},
  \texttt{random\_benchmark\_fields.npz};
\item \texttt{real\_porous\_benchmark\_3d/outputs/benchmark\_w1n/}:\\
  \texttt{benchmark\_metrics.csv}, \texttt{benchmark\_report.json};
\item the base-project verification run\\
  \texttt{ddpnm\_3d\_uniform\_spheres/outputs/verify\_run/} (Classic
  Schur eigenvalue footnote).
\end{itemize}
The geometry descriptors and tabulated numerical values are retained
from the supplied working experiment document.  In this revised version,
Figs.~\ref{fig:budget}, \ref{fig:metrics}, \ref{fig:pareto} and
\ref{fig:shares} were regenerated directly from the tabulated values;
the geometry and Random-27 field panels were reproduced from the supplied
compiled artifact.  No unreported experiment was assigned a numerical
value.

\end{document}
